\section{Conclusion}
We introduce Human and Mouse benchmarks that compare peptides from the same parent protein under the same MHC restriction and predict which has stronger presentation evidence in a specified tissue. The two members of every pair are additionally matched on their total number of positive tissue occurrences; a pair-level audit confirms zero mismatches and identical occurrence distributions between labels in both species and both partitions. The experiments use one fixed internal test set with 50 pairs per retained task and three common final-training seeds. Species-specific TissuePMHC configurations are selected by pair-grouped cross-validation confined to the training files; no peptide-separated generalization result is reported. Tuned TissuePMHC reaches mean AUROC/AUPRC values of 0.8136/0.8126 on Human and 0.8194/0.8279 on Mouse. Same-condition ablations support the multi-kernel encoder in both species and the MHC-specific branch in Human, whereas a trainable MHC-only control is worse than TissuePMHC in all 77 Human and all 11 Mouse tasks. Post-hoc branch-level SHAP analysis of the Human model localizes the strongest pair-matched sequence contributions to P2, P3, and P9 and recovers stable HLA-stratified residue patterns; exhaustive single-residue ISM independently recovers the same position ordering and strongly agrees with branch-level attribution. Holding peptide, HLA, and substitution fixed while switching tissue tasks produces greater perturbation dispersion than training-seed variation for all 28 multi-tissue HLA alleles, directly demonstrating fitted-model tissue conditioning. Primary tissue-independent external predictors also remain near random in both species. The retained signal concerns tissue-associated peptide ranking in previously represented tissue--MHC combinations; its SHAP and ISM structure describes model dependence and does not identify a specific antigen-processing mechanism. Validation in unseen peptides, tissues, alleles, studies, and prospective samples remains necessary before biological or clinical application.
